% TOP_PROBLEM: {"id": 5300036, "title": "Periodic point on a Siegel-disk boundary", "category_id": 11, "category": {"id": 11, "name": "dynamical_systems", "display_name": "Dynamical Systems", "description": "Problems about long-term behavior of deterministic systems, Hamiltonian dynamics, and periodic orbits.", "slug": "dynamical-systems", "order_index": 11, "created_at": "2026-07-31T15:26:25.671Z"}, "status": "partially_solved", "problem_number": "AMR-052-0036", "set_id": 13}
% FIRST_POSTED: 2026-10-02
% ENTRY_KIND: statement_only
% UnsolvedMath ID 5300036; source catalogue code AMR-052-0036
% !TeX program = lualatex
% Definition and sources only; this file is not an LLM solution attempt.
\documentclass[11pt,a4paper]{article}
\usepackage[margin=25mm]{geometry}
\usepackage{fontspec}
\setmainfont{lmroman10-regular.otf}[BoldFont=lmroman10-bold.otf,
 ItalicFont=lmroman10-italic.otf,BoldItalicFont=lmroman10-bolditalic.otf]
\usepackage{amsmath,amssymb,mathtools}
\usepackage{unicode-math}
\setmathfont{latinmodern-math.otf}
\AtBeginDocument{\let\setminus\smallsetminus}
\usepackage{xurl,hyperref}
\hypersetup{colorlinks=true,urlcolor=blue}
\setcounter{secnumdepth}{0}
\setlength{\parindent}{0pt}
\setlength{\parskip}{5pt}
\setlength{\emergencystretch}{3em}
\begin{document}
\section*{Periodic point on a Siegel-disk boundary}\label{5300036}
{\small UnsolvedMath ID 5300036 · AMR-052-0036 · Dynamical Systems · AMR Open Problem Lists}
\subsection{Definitions and mathematical statement}
For a rational map $R:\widehat{\mathbb C}\to\widehat{\mathbb C}$ of
degree at least two, the Fatou set is the largest open set on which
the iterates $\{R^n:n\ge0\}$ form a normal family in the spherical
metric. A Fatou component $U$ of least period $p\ge1$ is a Siegel
disk if there are a conformal bijection $h:\mathbb D\to U$ and an
irrational $\theta$ such that
\[
 R^p(h(z))=h(e^{2\pi i\theta}z)\qquad(z\in\mathbb D).
\]
Its boundary $\partial U$ is taken in the Riemann sphere.

The problem asks whether there can exist such an $R,U$ and a point
$a\in\partial U$ which is periodic for $R$:
\[
 R^q(a)=a\qquad\text{for some integer }q\ge1.
\]
The period $q$ need not initially be specified, and the map is allowed
to have any finite degree. A construction would have to prove both
that $U$ is the full Siegel Fatou component and that the periodic
point lies on its actual boundary.

The question does not concern the center $h(0)$, which is already a
periodic point inside $U$. It also does not assume that $h$ extends
continuously to the circle, that the boundary is a Jordan curve, or
that the rotation number has bounded continued-fraction entries.
Those regularity hypotheses can exclude the proposed phenomenon and
would restrict the general question. Accumulation of periodic points
near the boundary is different from a periodic point on the boundary.
\subsection{Short English statement}
Can the boundary of a rational map’s Siegel disk contain a periodic point?
\subsection{Sources}
\begin{itemize}
\item[S1] ULAM.AI and the UnsolvedMath Contributors, supplied problems.json, record 5300036 (AMR-052-0036), AMR Open Problem Lists. Catalogue identity and source transcription; CC BY 4.0.
\url{https://huggingface.co/datasets/ulamai/UnsolvedMath}.
\item[S2] Stony Brook - Open Problems in Dynamical Systems, 1992, Milnor local P11. Original source indexed by AMR-052-0036.
\url{https://www.math.stonybrook.edu/open-problems-dynamical-systems}.
\item[S3] Milnor, Problems on local connectivity, Problem11, in Bielefeld–Lyubich (1992).
\url{https://www.math.stonybrook.edu/preprints/ims92-7.pdf}.
\end{itemize}
\end{document}
